\clearpage
\begin{explanationbox}
\textbf{一、为什么会想到这种方法？}
\begin{center}\begin{tikzpicture}[font=\small,>=stealth,line join=round]
\draw[gray!40,rounded corners=3pt] (0,0) rectangle (14.2,8.3);\node at (7.1,7.700000000000001) {走一步变号、走两步还原：整个图像才重复};
\begin{scope}[shift={(1.7000000000000002,3.8062499999999995)},x=1.2cm,y=1.4374999999999998cm]\draw[->,gray] (-0.5,0)--(4.5,0) node[right]{$x$};\draw[->,gray] (0,-1.5)--(0,1.7) node[above]{$y$};\draw[blue,thick,solid,domain=-0.2:4.2,samples=180] plot(\x,{cos(deg(pi*\x))});\fill[orange!85!black] (0,1) circle(2.2pt);\fill[orange!85!black] (1,-1) circle(2.2pt);\fill[orange!85!black] (2,1) circle(2.2pt);\end{scope}\node[anchor=west,align=left,text width=5.0cm] at (8.4,4.0) {示例 $f(x)=\cos(\pi x)$\\[9pt]$f(0)=1$，$f(1)=-1$\\[9pt]$f(2)=1$\\[9pt]$f(x+1)=-f(x)$\\[9pt]$f(x+2)=f(x)$};
\node[text width=12.6cm,align=center] at (7.1,.65) {相隔1不保证同高；相隔2时，对每一个位置的高度都恢复。};
\end{tikzpicture}\end{center}



从最熟悉的正弦函数说起。$\sin x$ 的图像向右移 $2\pi$ 后，凹凸位置与高度全部重合。
这说明“周期”不是单独一个点回到了旧高度，而是\textbf{整幅图都在同样的距离后重复}。

现在题目不告诉我们图像，只有 $f(x+1)=-f(x)$。我们可以从任意 $x$ 出发，把 $f(x)$ 暂时记作 $2$：
\[
\begin{array}{c|ccccc}
 \text{位置}&x&x+1&x+2&x+3&x+4\\\hline
 \text{函数值}&2&-2&2&-2&2
\end{array}
\]
这仅是“假设起点取 $2$”时的演示。真正关键的是：\textbf{无论起点是什么数}，
负号的负号一定会还原。这才足以说明所有位置都重复。


\end{explanationbox}\clearpage\begin{explanationbox}\textbf{二、三种最容易混淆的情况}
\begin{center}\begin{tikzpicture}[font=\small,>=stealth,line join=round]
\draw[gray!40,rounded corners=3pt] (0,0) rectangle (14.2,8.3);\node at (7.1,7.700000000000001) {取倒数两次还原：两个高度的乘积始终为1};
\begin{scope}[shift={(1.4333333333333333,1.9907407407407407)},x=0.6666666666666666cm,y=1.7037037037037035cm]\draw[->,gray] (-0.5,0)--(8.5,0) node[right]{$x$};\draw[->,gray] (0,-0.2)--(0,2.5) node[above]{$y$};\draw[blue,thick,solid,domain=-0.2:8.2,samples=180] plot(\x,{pow(2,cos(deg(pi*\x/2)))});\fill[orange!85!black] (0,2) circle(2.2pt);\fill[orange!85!black] (2,0.5) circle(2.2pt);\fill[orange!85!black] (4,2) circle(2.2pt);\end{scope}\node[anchor=west,align=left,text width=5.0cm] at (8.4,4.0) {示例 $f=2^{\cos(\pi x/2)}$\\[9pt]$f(0)=2,\ f(2)=\frac12$\\[9pt]$f(4)=2$\\[9pt]$f(x+2)=\frac1{f(x)}$\\[9pt]处处为正，没有零值};
\node[text width=12.6cm,align=center] at (7.1,.65) {不是只让标出的三个点满足条件；由余弦变号可验证整条曲线的倒数关系。};
\end{tikzpicture}\end{center}


\begin{itemize}[leftmargin=*,itemsep=0.45em]
\item $f(x+2)=f(x)$：往右走 $2$，高度\textbf{已经回来}，所以 $2$ 是周期。
\item $f(x+2)=-f(x)$：往右走 $2$，高度变成相反数；再走 $2$ 才回来，所以 $4$ 是周期。
\item $f(x+2)=1/f(x)$：例如 $2\to1/2\to2$；换成任意非零数 $u$，也有
$u\to1/u\to u$，所以 $4$ 是周期。
\end{itemize}

\end{explanationbox}\clearpage\begin{explanationbox}\textbf{三、为什么“两个位置相等”要相减？}
\begin{center}\begin{tikzpicture}[font=\small,>=stealth,line join=round]
\draw[gray!40,rounded corners=3pt] (0,0) rectangle (14.2,8.3);\node at (7.1,7.700000000000001) {例3：比较实际输入之差，负号也要带进去};
\begin{scope}[shift={(1.85,3.8062499999999995)},x=0.75cm,y=1.4374999999999998cm]\draw[->,gray] (-1,0)--(7,0) node[right]{$x$};\draw[->,gray] (0,-1.5)--(0,1.7) node[above]{$y$};\draw[blue,thick,solid,domain=-0.5:6.5,samples=180] plot(\x,{cos(deg(2*pi*\x/3))});\fill[orange!85!black] (0,1) circle(2.2pt);\fill[orange!85!black] (3,1) circle(2.2pt);\end{scope}\node[anchor=west,align=left,text width=5.0cm] at (8.4,4.0) {示例 $f=\cos\frac{2\pi x}3$\\[9pt]$f(x+1)=f(x-2)$\\[9pt]输入差 $1-(-2)=3$\\[9pt]图中 $f(0)=f(3)=1$\\[9pt]3是一个正周期};
\node[text width=12.6cm,align=center] at (7.1,.65) {令 $t=x-2$ 后，两个输入成为 $t$ 与 $t+3$；不是比较常数的绝对值。};
\end{tikzpicture}\end{center}



给定 $f(x+1)=f(x-2)$。看两边的输入：$x+1$ 在右，$x-2$ 在左，相距
\[(x+1)-(x-2)=3.\]
也就是说，\textbf{相隔 $3$ 的两个位置，高度总一样}；所以周期是 $3$，不是 $1$ 或 $2$。
一般地，两边是 $x+a$ 与 $x+b$，就比较它们之间相隔多远：$|a-b|$。

\end{explanationbox}\clearpage\begin{explanationbox}\textbf{四、分式为什么也会两步还原？}
\begin{center}\begin{tikzpicture}[font=\small,>=stealth,line join=round]
\draw[gray!40,rounded corners=3pt] (0,0) rectangle (14.2,8.3);\node at (7.1,7.700000000000001) {分式证明一：输出再作为输入，回到原来的数};
\begin{scope}[shift={(2.975,3.7409090909090907)},x=0.75cm,y=0.41818181818181815cm]\draw[->,gray] (-2.5,0)--(5.5,0) node[right]{$u$};\draw[->,gray] (0,-5)--(0,6) node[above]{$v$};\draw[blue,thick,solid,domain=-2:0.6,samples=180] plot(\x,{(\x+1)/(\x-1)});\draw[blue,thick,solid,domain=1.5:5,samples=180] plot(\x,{(\x+1)/(\x-1)});\draw[gray,thick,dashed,domain=-2:5,samples=180] plot(\x,{\x});\fill[orange!85!black] (2,3) circle(2.2pt);\fill[orange!85!black] (3,2) circle(2.2pt);\draw[orange!85!black,->] (2,3)--(3,3)--(3,2)--(2,2);\end{scope}\node[anchor=west,align=left,text width=5.0cm] at (8.4,4.0) {$\Phi(u)=\frac{u+1}{u-1}$\\[9pt]图中 $2\to3\to2$\\[9pt]蓝曲线为 $v=\Phi(u)$\\[9pt]虚线为 $v=u$\\[9pt]一般需证 $\Phi(\Phi(u))=u$};
\node[text width=12.6cm,align=center] at (7.1,.65) {橙线把第一次输出移作第二次输入；字母通分证明覆盖所有 $u\ne1$。};
\end{tikzpicture}\end{center}


\begin{center}\begin{tikzpicture}[font=\small,>=stealth,line join=round]
\draw[gray!40,rounded corners=3pt] (0,0) rectangle (14.2,8.3);\node at (7.1,7.700000000000001) {分式证明二：把负分数送回同一规则也会还原};
\begin{scope}[shift={(3.725,3.9499999999999997)},x=0.75cm,y=0.3833333333333333cm]\draw[->,gray] (-3.5,0)--(4.5,0) node[right]{$u$};\draw[->,gray] (0,-6)--(0,6) node[above]{$v$};\draw[blue,thick,solid,domain=-3:-1.5,samples=180] plot(\x,{(1-\x)/(1+\x)});\draw[blue,thick,solid,domain=-0.7:4,samples=180] plot(\x,{(1-\x)/(1+\x)});\draw[gray,thick,dashed,domain=-3:4,samples=180] plot(\x,{\x});\fill[orange!85!black] (2,-0.3333333333333333) circle(2.2pt);\fill[orange!85!black] (-0.3333333333333333,2) circle(2.2pt);\draw[orange!85!black,->] (2,-.333333)--(-.333333,-.333333)--(-.333333,2)--(2,2);\end{scope}\node[anchor=west,align=left,text width=5.0cm] at (8.4,4.0) {$\Psi(u)=\frac{1-u}{1+u}$\\[9pt]图中 $2\to-\frac13\to2$\\[9pt]蓝曲线为 $v=\Psi(u)$\\[9pt]虚线为 $v=u$\\[9pt]一般需证 $\Psi(\Psi(u))=u$};
\node[text width=12.6cm,align=center] at (7.1,.65) {图中两次代入均合法；字母通分证明覆盖所有 $u\ne-1$。};
\end{tikzpicture}\end{center}



先别把它看成“函数”，把它当成一个对数做操作的小机器。
第一台机器收到 $u$，输出 $(u+1)/(u-1)$；第二台机器收到 $u$，输出 $(1-u)/(1+u)$。
\textbf{各自把自己的输出再放进去一次，就会得到原来的 $u$}。

这就解释了为什么它们与“变号”“取倒数”属于同一思路：
\textbf{只要走一步发生变化、走两步恢复原状，就能找到长度为两步的周期。}

\end{explanationbox}\clearpage\begin{explanationbox}\textbf{五、周期和最小正周期，差在哪里？}
\begin{center}\begin{tikzpicture}[font=\small,>=stealth,line join=round]
\draw[gray!40,rounded corners=3pt] (0,0) rectangle (14.2,8.3);\node at (7.1,7.700000000000001) {2是周期，却能容纳三次更短的完整重复};
\begin{scope}[shift={(1.3608695652173914,3.7409090909090907)},x=2.6086956521739126cm,y=1.3939393939393938cm]\draw[->,gray] (-0.1,0)--(2.2,0) node[right]{$x$};\draw[->,gray] (0,-1.5)--(0,1.8) node[above]{$y$};\draw[blue,thick,solid,domain=0:2,samples=180] plot(\x,{sin(deg(3*pi*\x))});\draw[<->,orange!85!black] (0,1.35)--(.6666667,1.35) node[midway,above]{$\frac23$};\end{scope}\node[anchor=west,align=left,text width=5.0cm] at (8.4,4.0) {$f=\sin(3\pi x)$\\[9pt]$f(x+1)=-f(x)$\\[9pt]两步给出周期2\\[9pt]最小正周期为 $\frac23$\\[9pt]$2=3\times\frac23$};
\node[text width=12.6cm,align=center] at (7.1,.65) {每个波峰间距为2/3；由条件推出的一个周期，未必是最小正周期。};
\end{tikzpicture}\end{center}



如果一张图每隔 $2$ 就重复，它每隔 $4$、$6$ 也重复。只知道周期为 $4$，不能说 $4$ 最小。
更意外的是：$f(x)=\sin(3\pi x)$ 的最小正周期为 $\frac23$，却满足
\[f(x+1)=-f(x).\]
因此由“走一步变号”得到的 $2$ 的确是一个周期，但离最小的 $\frac23$ 还很远。


\end{explanationbox}\clearpage\begin{explanationbox}\textbf{六、考场上什么时候用？}
\begin{center}\begin{tikzpicture}[font=\small,>=stealth,line join=round]
\draw[gray!40,rounded corners=3pt] (0,0) rectangle (14.2,8.3);\node at (7.1,7.700000000000001) {两次轴对称，把整幅图横移两轴间距的两倍};
\begin{scope}[shift={(1.9181818181818182,3.822222222222222)},x=0.5454545454545454cm,y=0.12777777777777777cm]\draw[->,gray] (-1.5,0)--(9.5,0) node[right]{$x$};\draw[->,gray] (0,-17)--(0,19) node[above]{$y$};\draw[blue,thick,solid,domain=-1:9,samples=180] plot(\x,{14*cos(deg(pi*(\x-1)/3))});\fill[orange!85!black] (2,7) circle(2.2pt);\fill[orange!85!black] (8,7) circle(2.2pt);\foreach \a in {1,4}{\draw[violet,dashed] (\a,-16)--(\a,16);\node[above,violet] at (\a,16){$x=\a$};}\end{scope}\node[anchor=west,align=left,text width=5.0cm] at (8.4,4.0) {示例 $f=14\cos\frac{\pi(x-1)}3$\\[9pt]两轴 $x=1,4$\\[9pt]两次配对 $x\to2-x\to x+6$\\[9pt]$f(2)=f(8)=7$\\[9pt]周期为6（只保证一个）};
\node[text width=12.6cm,align=center] at (7.1,.65) {本例曲线取振幅14以满足已知值7；正文推导适用于所有满足题设的函数。};
\end{tikzpicture}\end{center}


\begin{itemize}[leftmargin=*,itemsep=0.45em]
\item \textbf{求大数位置的值：}先找一个周期，再把 $2026$ 减去若干个周期，化为小数值位置。
\item \textbf{判断周期：}看函数值走一步如何变化，再计算走几步能回到原值。
\item \textbf{结合对称性：}两条不同的对称轴，也可能让图像经过两次翻折后平移出一个周期（第 03、04 部分有例子）。
\end{itemize}
\end{explanationbox}
